\documentclass{article}
\usepackage{amsmath, amssymb, amsthm}
\usepackage{hyperref}
\usepackage{geometry}
\geometry{margin=1in}

\title{Erd\H{o}s Problem \#131}
\date{Last updated: 2025-08-31}
\author{Source: erdosproblems.com}

\begin{document}
\maketitle

\noindent\textbf{Status:} OPEN \quad \textbf{No prize}

\noindent\textbf{Tags:} number theory

\medskip
\noindent\textbf{Problem Statement:}

Let $F(N)$ be the maximal size of $A\subseteq\{1,\ldots,N\}$ such that no $a\in A$ divides the sum of any distinct elements of $A\backslash\{a\}$. Estimate $F(N)$. In particular, is it true that\[F(N) > N^{1/2-o(1)}?\]

\bigskip
\noindent\textbf{Remarks:}

This was studied by Erd\H{o}s, Lev, Rauzy, S\'{a}ndor, and S\'{a}rk\"{o}zy \cite{ELRSS99}, where they call such a property 'non-dividing', and prove the explicit bound\[F(N)<3N^{1/2}+1.\]In \cite{Er97b} Erd\H{o}s credits Csaba with a construction that proves $F(N) \gg N^{1/5}$. Such a construction was also given in \cite{ELRSS99}, where it is linked to the problem of non-averaging sets (see [186]).

Indeed, every such set is non-averaging, and hence the result of Pham and Zakharov \cite{PhZa24} implies\[F(N) \leq N^{1/4+o(1)}.\]This shows the answer to the original question is no, but the general question of the correct growth of $F(N)$ remains open.

In \cite{Er75b} Erd\H{o}s writes that he originally thought $F(N) <(\log N)^{O(1)}$, but that Straus proved that\[F(N) > \exp((\sqrt{\tfrac{2}{\log 2}}+o(1))\sqrt{\log N}).\]See also [13].

This is discussed in problem C16 of Guy's collection \cite{Gu04}.

\bigskip
\noindent\textbf{References:}

[ELRSS99] Erd\H{o}s, P. and Lev, V. and Rauzy, G. and S\'{a}ndor, C. and
S\'{a}rk\"ozy, A., Greedy algorithm, arithmetic progressions, subset sums and
divisibility. Discrete Math. (1999), 119--135.
 
 [Er75b] Erd\H{o}s, Paul, Problems and results in combinatorial number theory. Journ\'{e}es Arithm\'{e}tiques de Bordeaux (Conf., Univ. Bordeaux, Bordeaux, 1974) (1975), 295-310.
 
 [Er97b] Erd\H{o}s, Paul, Some old and new problems in various branches of combinatorics. Discrete Math. (1997), 227-231.
 
 [Gu04] Guy, Richard K., Unsolved problems in number theory. (2004), xviii+437.
 
 [PhZa24] Pham, H. T. and Zakharov, D., Sharp bound for the Erd\H{o}s-Straus non-averaging set problem. arXiv:2410.14624 (2024).

\bigskip
\noindent\small{Source: \url{https://www.erdosproblems.com/131}}
\end{document}
